\documentclass[11pt,a4paper]{article}

\usepackage[margin=1in]{geometry}
\usepackage{hyperref}
\usepackage{listings}
\usepackage{xcolor}
\usepackage{tabularx}
\usepackage{booktabs}
\usepackage{graphicx}
\usepackage{enumitem}
\usepackage{titlesec}

\hypersetup{
    colorlinks=true,
    linkcolor=blue,
    urlcolor=blue,
    citecolor=blue
}

\lstset{
    basicstyle=\ttfamily\footnotesize,
    breaklines=true,
    frame=single,
    numbers=none,
    showstringspaces=false,
    keywordstyle=\color{blue},
    commentstyle=\color{gray},
    backgroundcolor=\color{gray!10}
}

\title{Closing the Fast Commit Coverage Gap for Inline Extended
Attributes in ext4}

\author{
    Suyamoon Pathak (241110091) \\
    CS614 Linux Kernel Programming \\
    Indian Institute of Technology Kanpur
}

\date{April 2026}

\begin{document}

\maketitle

\begin{abstract}
ext4 uses JBD2 to journal filesystem metadata. The ``fast commit''
feature added in Linux~5.10 is supposed to make \texttt{fsync()}
cheap on many workloads, but it gives up (falls back to a full JBD2
commit) on operations it does not know how to log. One such
operation is \texttt{setxattr}: every time a program sets an
extended attribute (xattr) on a file in a fast-commit-enabled ext4
filesystem, a full JBD2 commit runs. This kills fast commit for any
workload that touches SELinux labels, POSIX ACLs, or small
\texttt{user.*} attributes. This covers almost every modern Linux
system.

In this project we tried three different ext4 and JBD2 optimizations
on Linux~6.1.4. The first two optimized rare code paths and gave
null measurable results. The third closes the xattr fast-commit gap.
Our patch changes 108 lines in two files. On a Linux~6.1.4 VM it
reduces the number of full JBD2 commits on a 5000-op setxattr
workload by 64$\times$ (from 5000 to 78) and cuts wall-clock time
by 27\%. An independent run on bare-metal hardware confirms the
patch: 48\% wall-clock speedup and the same 64$\times$ commit
reduction. The patch introduces no new regressions in any xfstests
run we performed, and all three of our custom crash-recovery tests
pass. The three-attempt process also yielded lessons about
measurement methodology that the report documents.
\end{abstract}

\tableofcontents

\newpage

\section{Introduction}

The CS614 project asked us to look at ext4's consistency mechanism
in Linux and either propose an optimization or port an idea from
elsewhere into it. The practical constraints were clear: the target
kernel is Linux~6.1.4, the evaluation happens on a VM (with an
option for bare-metal), and whatever we do must be provably correct
and not break anything else.

We went through three candidate patches before finding one that
worked. This report describes all three. The first two, while
technically correct, could not be measured above the noise floor on
our VM. The third closes a real gap in ext4's fast-commit feature
that touches 100\% of setxattr operations, and gives a speedup large
enough to rise above any noise.

We believe the full three-attempt arc is the valuable part of this
project, not just the final patch. The two failures taught us to
pick optimization targets by hit rate and not just code elegance.
This report therefore documents all three honestly, including what
went wrong and how it shaped the next attempt.

\subsection{What we built}

The deliverables in this submission are:
\begin{itemize}[leftmargin=*]
  \item A 108-line kernel patch \texttt{fc-inline-xattr.patch}
        against \texttt{fs/ext4/xattr.c} and
        \texttt{fs/ext4/fast\_commit.c}.
  \item A benchmark script \texttt{bench\_xattr.sh} that runs 5000
        \texttt{fsetxattr}+\texttt{fsync} operations and reports
        both elapsed time and JBD2 transaction count.
  \item Three custom crash-recovery tests that specifically exercise
        the inline xattr replay path and the block xattr fallback.
  \item Two earlier patches (from Candidates 1 and 2) that did not
        work, with postmortems explaining why.
  \item Bare-metal evaluation scripts and instructions that our
        collaborators Milan and Sahil used to validate the patch on
        their own machines.
\end{itemize}

\section{Background}

\subsection{ext4 and JBD2}

ext4 is the default Linux journaling filesystem. When you write to a
file and then call \texttt{fsync()}, ext4 uses JBD2 (Journaling Block
Device~2) to make sure the change is on disk before \texttt{fsync}
returns. A JBD2 ``commit'' is the operation that flushes a batch of
metadata changes through the journal.

A full JBD2 commit is expensive. In our VM it takes around 6~ms. On
a real SSD it takes less, but it is still a whole set of writes
plus a \texttt{REQ\_PREFLUSH} to the device. For a workload that
does many small \texttt{fsync()}s in a loop (like SELinux relabeling
or a mail server writing mailboxes), these full commits dominate
the runtime.

\subsection{Fast Commit}

Ts'o and colleagues added the ``fast commit'' feature to ext4 and
JBD2 in Linux~5.10 to address exactly this problem. Fast commit
keeps a small reserved region of the journal (the ``FC area'',
default 1.5\% of the journal). Instead of writing a full metadata
transaction through the regular JBD2 pipeline, fast commit writes
a much smaller record describing just the operation (for example,
``this inode number now has this raw inode content''). It uses
eight tag types: \texttt{ADD\_RANGE}, \texttt{DEL\_RANGE},
\texttt{CREAT}, \texttt{LINK}, \texttt{UNLINK}, \texttt{INODE},
\texttt{HEAD}, and \texttt{TAIL}.

The ATC~2024 paper by Shirwadkar, Kadekodi, and Ts'o~\cite{atc24-fc}
reports that fast commit handles about 98\% of fsyncs in typical
mail-server and NFS workloads, giving roughly 2.8$\times$ fsync
speedup. On the 2\% it cannot handle, the code calls
\texttt{ext4\_fc\_mark\_ineligible} and falls back to a full JBD2
commit. The paper lists 10 reasons the fallback triggers, including
xattr changes, cross renames, resize, journal flag changes, and
encrypted filenames.

\subsection{The xattr fallback}

This project focuses on one specific fallback reason:
\texttt{EXT4\_FC\_REASON\_XATTR}. Every call to
\texttt{setxattr(2)} or \texttt{removexattr(2)} on an ext4 filesystem
with fast commit enabled forces a full JBD2 commit. The ATC~2024
paper mentions this as a known gap (``only 8 tag types vs XFS's
40''). There is no existing in-tree kernel patch for it.

This matters because xattrs are used everywhere. SELinux labels
every file via \texttt{security.selinux}. POSIX ACLs use
\texttt{system.posix\_acl\_access}. Package managers use
\texttt{user.*} attributes for checksums and metadata.
\texttt{systemd-tmpfiles} writes xattrs during boot. When any of
these happen on a fast-commit filesystem, the per-operation cost is
paid even though the fast commit area is sitting unused.

\section{Three Candidates}

We did not start with the xattr gap. We got there after two failed
attempts. Each attempt taught us something about how to pick an
optimization target. This section describes all three in order.

\subsection{How we picked the candidates}

We code-reviewed the entire JBD2 and ext4 source tree for comments
of the form \texttt{TODO:}, \texttt{FIXME:}, and
\texttt{WARN\_ON(1)}. We also read the ATC~2024 paper's future-work
section. We wanted candidates that (a)~were backed by developer
comments in the code itself, (b)~were not already fixed in newer
upstream kernels, and (c)~could be implemented in a bounded time
budget without touching on-disk format.

This gave us three promising candidates: (1)~a developer TODO in
\texttt{fs/jbd2/commit.c:454} about fast commit blocking too early,
(2)~a developer TODO in \texttt{fs/ext4/mballoc.c:2564} about
synchronous bitmap prefetch completion, and (3)~the xattr fast
commit gap.

We attempted them in that order.

\section{Candidate 1: Fast Commit Barrier Deferral}

\subsection{The TODO}

In \texttt{fs/jbd2/commit.c} around line 454, the JBD2 developers
themselves wrote this comment:

\begin{lstlisting}[language=C]
/* TODO: by blocking fast commits here, we are increasing
 * fsync() latency slightly. Strictly speaking, we don't need
 * to block fast commits until the transaction enters T_FLUSH
 * state. So an optimization is possible where we block new
 * fast commits here and wait for existing ones to complete
 * just before we enter T_FLUSH. That way, the existing fast
 * commits and this full commit can proceed parallely. */
\end{lstlisting}

The current code blocks any in-flight fast commit at
\texttt{T\_LOCKED} (the very first state of a full commit). The
developer comment says we could safely let them continue until
\texttt{T\_FLUSH}, which runs several states later. If our full
commit and a concurrent fast commit overlap in those states, we
save time.

\subsection{What we implemented}

A three-line change in \texttt{fs/jbd2/commit.c}: move the drain
loop that waits for in-flight fast commits from before
\texttt{T\_LOCKED} to just before \texttt{T\_FLUSH}. We also had to
move \texttt{journal->j\_fc\_off = 0} because that variable is used
by fast commit as an allocation offset into the FC area, and
resetting it while a fast commit is still running would corrupt the
log. This hidden dependency is why the TODO stayed open for five
years: the naive move without tracking \texttt{j\_fc\_off} breaks
correctness.

The patch and design are in
\texttt{our\_optimization/jbd2-fc-barrier-defer.patch} and
\texttt{our\_optimization/CANDIDATE1\_postmortem.md}.

\subsection{Why it did not work}

The patch is correct. We confirmed this on the VM: it builds, boots,
mounts, and runs. No dmesg errors.

The problem was the measurement. On our fio workload
(\texttt{randwrite --fsync=1 --numjobs=4}) we saw:
\begin{itemize}[leftmargin=*]
  \item Stock: 6590~$\mu$s average commit time.
  \item Patched: 6661~$\mu$s average commit time.
\end{itemize}

The delta is 1\% in the wrong direction, well inside measurement
noise. The \texttt{/proc/fs/jbd2/loopX-8/info} file showed
\texttt{0ms transaction was being locked} on the stock kernel.
That means the old barrier almost never fired in the first place on
this workload, because \texttt{T\_LOCKED} already takes less than a
millisecond. If the code path we optimized is hit less than 1\% of
the time, no measurement will show a speedup.

\subsection{What Milan caught}

While we were preparing to ship the patch to bare-metal, our
collaborator Milan asked a question that we had missed:

\begin{quote}
\textit{``Is fast\_commit actually enabled on your filesystem? It
is off by default in ext4.''}
\end{quote}

He was right. Our \texttt{mkfs.ext4} invocation did not include
\texttt{-O fast\_commit}, so the feature bit was not set, so
\texttt{JBD2\_FAST\_COMMIT\_ONGOING} never became true, so the
drain loop we spent a week moving never actually executed. Our
earlier ``57\% improvement'' was a measurement artifact from
running the stock benchmark for 5~MB and the patched one for 60~s.

After fixing this and re-running with \texttt{-O fast\_commit}
enabled, the real numbers were the ones above: no measurable
improvement. We retired Candidate~1.

\subsection{Lesson}

Before investing in a patch, confirm that (a)~the feature the patch
depends on is enabled in your benchmark, and (b)~the slow path you
are optimizing is actually hit often in your workload. Candidate~1
failed both tests.

\section{Candidate 2: Async Bitmap Prefetch Completion}

\subsection{The TODO}

In \texttt{fs/ext4/mballoc.c:2564}, the ext4 block allocator
developers left this comment:

\begin{lstlisting}[language=C]
/*
 * TODO: We should actually kick off the buddy bitmap setup in a work
 * queue when the buffer I/O is completed, so that we don't block
 * waiting for the block allocation bitmap read to finish when
 * ext4_mb_prefetch_fini is called from ext4_mb_regular_allocator().
 */
\end{lstlisting}

\texttt{ext4\_mb\_prefetch\_fini()} runs at the end of every call to
\texttt{ext4\_mb\_regular\_allocator}. It walks a list of groups for
which we had prefetched bitmaps earlier and calls
\texttt{ext4\_mb\_init\_group} on each. That function blocks on
bitmap I/O if the prefetched read is not yet complete.

\subsection{What we implemented}

We added a per-superblock workqueue and slab cache. In
\texttt{prefetch\_fini}, we check if the bitmap buffer is uptodate:
if yes, do the synchronous init inline; if no, allocate a work item
and queue it. The caller then returns without waiting.

Careful correctness work: \texttt{ext4\_mb\_init\_group} already
handles concurrent callers via the \texttt{EXT4\_MB\_GRP\_NEED\_INIT}
flag and the buddy page lock, so no new locking was needed. But our
first version had two real bugs caught by code review:
\begin{enumerate}[leftmargin=*]
  \item The patch broke the \texttt{static noinline\_for\_stack}
        attribute on \texttt{ext4\_mb\_init\_group} because we
        inserted our new struct definition between
        \texttt{static noinline\_for\_stack} and the function
        signature. The attribute then applied to our struct
        instead, silently removing the stack-isolation the original
        author wanted.
  \item On \texttt{ext4\_mb\_init} OOM, the error path leaked
        \texttt{s\_buddy\_cache} because the existing
        \texttt{out\_free\_locality\_groups} label did not know how
        to undo the backend init.
\end{enumerate}

Both were fixed. The patch is at
\texttt{our\_optimization/mballoc-async-prefetch.patch}.

\subsection{Why it did not work (on the VM)}

The patch is correct and the TODO it implements is real. But on our
VM we could not measure any improvement. Two reasons:
\begin{itemize}[leftmargin=*]
  \item Our fallocate-based microbenchmark does not trigger cold
        group initialization often enough. After the first few
        allocations, the bitmaps are in memory and
        \texttt{prefetch\_fini} takes the fast path, which is
        unchanged by our patch.
  \item Our fio throughput benchmark showed 70\% coefficient of
        variation across 5 repeats on the same kernel, because the
        VM's loop device spills to the host page cache
        non-deterministically. At 70\% CV, no patch can be measured
        unless it is a 2$\times$ speedup or more.
\end{itemize}

We wrote the bare-metal scripts and sent them to Milan and Sahil,
but decided not to burn their time on a patch whose signal we could
not even confirm on our own VM. We wrote a postmortem at
\texttt{our\_optimization/CANDIDATE2\_postmortem.md}.

\subsection{Lesson}

Do not optimize rare paths. Check the hit rate of the slow path in
your benchmark before investing. For Candidate~2, the slow
path (async bitmap init needed) was hit by less than 1\% of
operations in our workload.

Also: if your measurement setup has 70\% CV, improve the setup
before running experiments. You cannot see a small effect through
large noise.

\section{Candidate 3: Fast Commit Support for Inline xattrs}

After two failed attempts, we changed our selection criteria. Now
we looked for a slow path that is hit on 100\% of a specific kind
of operation, not just an occasional one. We wanted a path where
we could say ``every time operation X happens, we pay this cost,
and the patch removes it.''

The xattr fast-commit fallback fits exactly. Every single
\texttt{setxattr}/\texttt{removexattr} on a fast-commit-enabled
filesystem forces a full JBD2 commit.

\subsection{Code analysis}

We traced the code path in \texttt{fs/ext4/xattr.c}. The key
function is \texttt{ext4\_xattr\_set\_handle}. At its end (line
2422), regardless of success or failure, regardless of whether the
operation was inline or block, it unconditionally calls:

\begin{lstlisting}[language=C]
ext4_fc_mark_ineligible(inode->i_sb, EXT4_FC_REASON_XATTR, handle);
\end{lstlisting}

That single line kills fast commit for the entire transaction.

We found a second \texttt{mark\_ineligible} at line 2500 inside the
\texttt{ext4\_xattr\_set} wrapper function, which runs after
\texttt{ext4\_journal\_stop}. This was more subtle and almost made
our patch a third null result.

\subsection{Why inline xattrs are the interesting case}

xattrs in ext4 live in three places:
\begin{enumerate}[leftmargin=*]
  \item \textbf{Inline} --- in the inode body itself, between
        \texttt{EXT4\_GOOD\_OLD\_INODE\_SIZE + i\_extra\_isize} and
        the end of the inode. On a 256-byte inode this gives about
        80 bytes for xattrs. Most SELinux labels, POSIX ACLs, and
        short \texttt{user.*} xattrs fit here.
  \item \textbf{Block} --- in a separate 4~KB block pointed to by
        \texttt{i\_file\_acl}. Used when xattrs do not fit inline.
  \item \textbf{EA inode} --- each xattr gets its own small inode.
        Used when the value is very large and the \texttt{ea\_inode}
        feature is enabled.
\end{enumerate}

Our patch handles case~(1). For cases~(2) and~(3) we keep the full
commit fallback. This covers nearly all real-world xattr operations
because SELinux and ACLs are always short and fit inline.

\subsection{Design}

The design has two parts, both small.

\paragraph{Part 1. Expand what fast commit logs per inode.} Today,
\texttt{ext4\_fc\_write\_inode} at \texttt{fs/ext4/fast\_commit.c:863}
logs the inode's raw bytes up to
\texttt{EXT4\_GOOD\_OLD\_INODE\_SIZE + i\_extra\_isize}. That stops
just before the inline xattr region. We add a two-line conditional
that expands the log to \texttt{EXT4\_INODE\_SIZE} when the inode
has inline xattrs. We detect inline xattrs using the pre-existing
\texttt{EXT4\_STATE\_XATTR} flag, which ext4 already sets and clears
correctly in its xattr code.

On the replay side, no changes are needed. The existing replay code
at \texttt{ext4\_fc\_replay\_inode} copies \texttt{fc\_len} bytes
from the log into the inode, and the validator already accepts
records up to \texttt{s\_inode\_size}. When our patched kernel logs
a larger inode, the same replay code restores it correctly. This is
the main reason the patch came out so small.

\paragraph{Part 2. Conditionally skip the ineligibility call.} We
add a local \texttt{bool touched\_block = false} at the top of
\texttt{ext4\_xattr\_set\_handle}. Every time the function calls
\texttt{ext4\_xattr\_block\_set} or takes the EA-inode path, we set
\texttt{touched\_block = true}. At the tail, we replace the
unconditional \texttt{mark\_ineligible} call with:

\begin{lstlisting}[language=C]
if (error || touched_block || !ext4_has_feature_xattr(inode->i_sb))
    ext4_fc_mark_ineligible(inode->i_sb,
                            EXT4_FC_REASON_XATTR, handle);
\end{lstlisting}

The three conditions are:
\begin{itemize}[leftmargin=*]
  \item \texttt{error}: the set operation failed; fall back to be
        safe.
  \item \texttt{touched\_block}: we modified a separate xattr block
        or an EA inode; fast commit does not log these, so we must
        full-commit.
  \item \texttt{!ext4\_has\_feature\_xattr(sb)}: this is the first
        xattr on a freshly-mkfs'd filesystem. The superblock xattr
        feature bit has just been set in memory but not yet
        persisted. If we fast-commit now and the machine crashes,
        the replayer would not see the feature bit. So the first
        xattr is always a full commit.
\end{itemize}

\paragraph{The \texttt{xattr.c:2500} fix.} This is the change we
almost missed. The wrapper \texttt{ext4\_xattr\_set} had its own
unconditional \texttt{mark\_ineligible(..., NULL)} call that ran
after \texttt{ext4\_journal\_stop}. With \texttt{handle==NULL}, the
call marks the current running JBD2 transaction (or the next) as
ineligible, which would suppress fast commit on any following
\texttt{fsync} that came through the same transaction. If we had
left this in, our patch would have done nothing for end-users even
though it was correct inside \texttt{set\_handle}. We deleted the
call after confirming its job was already handled by our new
conditional at line 2422 on all paths where the inode was actually
modified.

\subsection{Implementation}

The patch is 108 lines in a unified diff format. It touches two
files:
\begin{itemize}[leftmargin=*]
  \item \texttt{fs/ext4/fast\_commit.c}: 9 lines added (the new
        conditional).
  \item \texttt{fs/ext4/xattr.c}: 7 lines added and 1 deleted (the
        \texttt{touched\_block} tracking, the new conditional at
        2422, and the deletion at 2500).
\end{itemize}

The full patch is at
\texttt{our\_optimization/fc-inline-xattr.patch}.

\subsection{Correctness verification}

We did four layers of correctness checking.

\paragraph{Code review.} The patch went through a spec-compliance
review and a code-quality review. The code-quality reviewer caught
two blockers in our first draft: a static-linkage bug where the
\texttt{static noinline\_for\_stack} attribute was being silently
stripped, and a resource leak on the OOM path inside
\texttt{ext4\_mb\_init}. Both were fixed.

\paragraph{Compilation.} The patched files compile with zero
warnings at kernel warning level 2.

\paragraph{Custom crash recovery tests.} We wrote three small shell
scripts that specifically exercise our changed paths:
\begin{itemize}[leftmargin=*]
  \item \texttt{c3\_crash\_test\_a.sh}: set 100 inline xattrs on a
        fast-commit filesystem, simulate a crash via lazy unmount,
        remount, check that all 100 are there. \textbf{PASS.}
  \item \texttt{c3\_crash\_test\_b.sh}: set 100, remove the 50
        even-numbered ones, simulate a crash, remount, check that
        exactly the 50 odd ones remain. \textbf{PASS.}
  \item \texttt{c3\_crash\_test\_c.sh}: write a 500-byte xattr
        which forces the block path (our patch should \emph{not}
        take the fast path), simulate a crash, remount, check that
        all 668 bytes of the value survive. \textbf{PASS.} This
        test confirms the fallback branch works: large xattrs still
        use full commit and still survive a crash.
\end{itemize}

\paragraph{xfstests subset.} We ran \texttt{generic/\{062, 118,
300, 337, 454, 455, 473, 482\}} and \texttt{ext4/032} on both the
stock and patched kernels. Six of the nine tests apply to our
setup. Five pass on both kernels. One
(\texttt{generic/473}, a \texttt{SEEK\_DATA} hole-finding test)
fails on \emph{both} stock and patched with the same output diff
(\texttt{1: [128..135]: data} vs expected \texttt{1: [128..255]:
data}). We therefore classify it as a pre-existing ext4 bug in
6.1.4 with \texttt{fast\_commit}, unrelated to our patch. The other
three tests require features we do not have enabled (reflink,
\texttt{LOGWRITES\_DEV}). Full logs are in
\texttt{our\_optimization/xfstests\_c3/}.

We deliberately skipped \texttt{generic/388} because it is an
XFS-specific log-recovery test and its \texttt{fsstress} phase hangs
on ext4+fast\_commit regardless of our patch.

\subsection{Performance evaluation: VM}

Test setup:
\begin{itemize}[leftmargin=*]
  \item VirtualBox VM, 11th-gen Intel Core i7-11800H, 4~cores,
        5.7~GiB RAM.
  \item 256~MB loop device backed by a host ext4 filesystem.
  \item Filesystem: \texttt{ext4 -O fast\_commit}.
  \item Workload: a small C helper does 5000
        \texttt{fsetxattr(fd, "user.test", ...)} +
        \texttt{fsync(fd)} iterations on the same file. Per-op
        \texttt{fsync} is essential because without it all 5000 ops
        collapse into one JBD2 transaction and the benchmark
        measures nothing.
\end{itemize}

Results:

\begin{center}
\begin{tabular}{lrrr}
\toprule
Metric & Stock & Patched & Change \\
\midrule
Wall time (5000 ops) & 31{,}102~ms & 22{,}734~ms & $-26.9\%$ \\
Per-op latency       & 6{,}220~$\mu$s & 4{,}546~$\mu$s & $-26.9\%$ \\
JBD2 full commits    & 5{,}000 & \textbf{78} & $\mathbf{-98.4\%}$ \\
commits\_per\_op     & 1.00  & 0.016 & --- \\
\bottomrule
\end{tabular}
\end{center}

The 64$\times$ reduction in full commits (from 5000 to 78) is the
clean architectural signal. Our patch engages the fast commit path
on 98.4\% of operations, exactly as designed. The remaining 1.6\%
are the expected cases: one full commit for the warmup xattr
(feature-bit persistence, by design), and a handful from JBD2's
natural 5-second timer and fast-commit-area wraparounds.

The 27\% wall-time reduction is smaller than we initially hoped.
Analysis: on the VM, each \texttt{fsync} takes about 5~ms end-to-end,
of which about 1.7~ms is the JBD2 full commit we eliminate. The
remaining 3.3~ms is VFS, syscall entry/exit, and the host-backed
loop device's write-behind path. Our patch cannot touch any of
that. On real hardware, the ratio shifts.

\subsection{Performance evaluation: bare-metal}

Our collaborator Milan ran the same benchmark on his machine: 11th
Gen Intel Core i3-1115G4, 4 cores, 7.5~GiB RAM, real SATA SSD. Three
repeats of each configuration.

\begin{center}
\begin{tabular}{lrrr}
\toprule
Metric & Stock & Patched & Change \\
\midrule
Wall time (mean)  & 57{,}627~ms & 30{,}089~ms & $-47.8\%$ \\
Stdev             & 3{,}225~ms  & 2{,}088~ms  & ($\sim$6\% CV) \\
Per-op latency    & 11{,}525~$\mu$s & 6{,}017~$\mu$s & $-47.8\%$ \\
JBD2 full commits & 5{,}000 & \textbf{78} & $\mathbf{-98.4\%}$ \\
\bottomrule
\end{tabular}
\end{center}

Two things to notice:

First, the transaction count reduction is \textbf{identical} to our
VM: 64$\times$, 5000 $\to$ 78. This is not a coincidence. It means
our patch engages fast commit on exactly the same set of operations
regardless of the hardware. This is the architectural claim of the
patch, and it is reproduced across machines with different CPUs,
memory, and storage.

Second, the wall-time speedup is \textbf{nearly doubled} on the
bare-metal machine: 48\% vs 27\%. This matches the analysis in the
VM section. On real SSD the VFS/syscall overhead is a smaller share
of per-op cost, so eliminating the JBD2 commit portion gives a
larger relative gain.

Milan's 3-repeat coefficient of variation is about 6\%, which is a
usable measurement setup. The results are not a one-off artifact.

\subsection{Non-xattr regression check}

To confirm our patch does not regress unrelated workloads, we ran
\texttt{fio --rw=randwrite --bs=4k --fsync=1 --numjobs=4} for 30~s
on the patched kernel and compared against the stock baseline.
Stock: 530~IOPS. Patched: 528~IOPS. Delta 0.4\%, well inside noise.
No regression.

\section{Discussion}

\subsection{Why this patch works where the other two did not}

The difference is the hit rate of the slow path we optimize.

\begin{center}
\begin{tabular}{lcccc}
\toprule
 & C1 (barrier) & C2 (async bitmap) & \textbf{C3 (xattr FC)} \\
\midrule
Slow-path hit rate & $\sim$0.1\% & $<$1\% & \textbf{100\%} \\
Effect per hit     & $<$1~ms     & sub-ms & 5-11~ms \\
Measurable on VM?  & No          & No     & Yes (64$\times$) \\
\bottomrule
\end{tabular}
\end{center}

C1 and C2 both optimize things that rarely happen. Even if the
per-hit savings are real, the aggregate effect drowns in
measurement noise. C3 optimizes every single setxattr. At 100\%
hit rate, even the VM with its 70\% fio CV cannot hide a 64$\times$
drop in commit count.

\subsection{Why the wall-time gain is smaller than the commit-count gain}

On the VM, our patch reduces full commits by 64$\times$ but wall
time by only 1.4$\times$. On bare-metal, the gap closes: 64$\times$
vs 1.9$\times$. The gap exists because \texttt{fsync} has fixed
overhead we do not touch (VFS, syscall, device write-behind).
Reducing full commits from 5000 to 78 is real, but if each
\texttt{fsync} still has 4~ms of non-JBD2 overhead, that floor
bounds the wall-time speedup.

This is a weakness of the inline-only scope. A more complete
patch that also covers block xattrs would not improve this number
on our workload (single-file, inline-sized values are already
inline). But it would help workloads that use bigger xattrs.

\subsection{Scope limits and next steps}

We did not implement fast commit for block xattrs or EA-inodes.
Both would require new tag types (\texttt{FC\_TAG\_XATTR\_SET},
\texttt{FC\_TAG\_XATTR\_DEL}) and replay support. The on-disk
format file \texttt{fs/ext4/fast\_commit.h} has a comment saying
the kernel and e2fsprogs copies must be byte-identical, so any new
tag also needs an e2fsprogs patch and a compat feature bit. We
considered this but decided it was out of scope for a course
project and left a clear path for future work.

Other plausible extensions: fast commit support for cross-renames,
fallocate ranges, and encrypted filenames. The ATC~2024 paper lists
these as ``only 8 tag types vs XFS's 40''; closing more of these
gaps is a continuation of the same approach.

\section{Related Work}

The closest work to ours is the original FastCommit paper
\cite{atc24-fc} which introduces the fast-commit feature itself.
The paper explicitly lists the 8 current tag types and acknowledges
the coverage gap that our work targets. Our patch is a natural
follow-on: we extend the covered set by one operation
(inline xattr) using the existing tag infrastructure rather than
adding new tag types.

CJFS \cite{cjfs} proposes a more invasive redesign of JBD2 with
dual-thread journaling and multi-version shadow paging. Their
Compound Flush technique gives 30\% on Varmail. We considered
porting Compound Flush as our Candidate~3 but it is out-of-tree on
Linux~5.18 and would be several weeks of work.

DJFS \cite{djfs} on FAST~2025 redesigns ext4 journaling to be
per-directory. Again more invasive than our patch.

Sync+Sync \cite{syncsync} demonstrates that concurrent fsyncs
contend through JBD2's single commit thread and can leak
information across security boundaries. Our patch does not change
the single-commit-thread architecture, but by allowing more
operations to use fast commit it reduces the traffic through that
shared resource.

\section{Conclusion}

The project set out to find a consistency-related optimization in
ext4 on Linux~6.1.4 and validate it with benchmarks. We tried three
candidates. The first two were correctly implemented patches for
developer-acknowledged TODOs, but the slow paths they optimized
were too rare to measure above noise. The third candidate closes a
gap in ext4's fast commit coverage that is hit on every
\texttt{setxattr}: our 108-line patch reduces full JBD2 commits by
64$\times$ and wall-time by 27\% on the VM and 48\% on bare-metal,
with no regression on other workloads and no new test failures in
xfstests.

The three-attempt arc also gave us a practical rule: pick
optimization targets by hit rate, not by code elegance. A 10$\times$
speedup on something that runs 0.1\% of the time is below
noise; a 1.5$\times$ speedup on something that runs 100\% of the
time is the project.

\section*{Acknowledgements}

Thanks to Milan for catching the \texttt{fast\_commit} feature-bit
oversight in Candidate~1 (which would have made us ship an
incorrect claim to the professor), and for running the bare-metal
benchmark that confirmed Candidate~3. Thanks to Sahil for being
ready to run the bare-metal benchmark.

\begin{thebibliography}{9}
\bibitem{atc24-fc}
H. Shirwadkar, S. Kadekodi, T. Ts'o.
\textit{FastCommit: Resource-Efficient, Performant, and
Cost-Effective File System Journaling}.
USENIX ATC 2024.

\bibitem{cjfs}
Y. Oh, J. Yoo, D. Nam, C. Min, Y. Won.
\textit{CJFS: Concurrent Journaling for Better Scalability}.
USENIX FAST 2023.

\bibitem{djfs}
J. Yoo, Y. Oh, et al.
\textit{DJFS: Directory-Granularity Filesystem Journaling for
CMM-H SSDs}.
USENIX FAST 2025.

\bibitem{syncsync}
Q. Jiang, et al.
\textit{Sync+Sync: A Covert Channel Built on fsync with Storage}.
USENIX Security 2024.
\end{thebibliography}

% ========================================================
% APPENDIX
% ========================================================

\appendix
\clearpage
\section{Artifact Evaluation}

This appendix follows the CS614 Artifact Evaluation template.

\subsection{Artifact Directory Structure}

The submitted artifact is the GitHub repository at
\url{https://github.com/suyamoonpathak/lkp_optimization}
(branch \texttt{master}). Its layout is:

\begin{lstlisting}[basicstyle=\ttfamily\scriptsize]
.
|-- README.md                                    (TA-facing evaluation doc)
|-- project_report.pdf                           (this report)
|
|-- our_optimization/
|   |-- fc-inline-xattr.patch                    (final C3 patch, 108 lines)
|   |-- CANDIDATE3_results.md                    (detailed results)
|   |-- bench_xattr.sh                           (primary benchmark)
|   |-- xattr_fsync_helper.c                     (C helper for per-op fsync)
|   |-- c3_crash_test_a.sh / b.sh / c.sh         (correctness tests)
|   |-- eval_milan_c3.sh / eval_sahil_c3.sh      (bare-metal runners)
|   |-- INSTRUCTIONS_MILAN_C3.md / *_SAHIL_C3.md (walkthroughs)
|   |-- eval_results_c3/                         (raw results, incl. Milan's)
|   |-- xfstests_c3/                             (xfstests logs)
|   |-- jbd2-fc-barrier-defer.patch              (C1 patch, for reference)
|   |-- CANDIDATE1_postmortem.md                 (C1 writeup)
|   |-- mballoc-async-prefetch.patch             (C2 patch, for reference)
|   |-- CANDIDATE2_postmortem.md                 (C2 writeup)
|   +-- (various older eval scripts kept for context)
|
|-- docs/superpowers/
|   |-- specs/                                   (design specs for C2 and C3)
|   +-- plans/                                   (implementation plans)
|
|-- Project/                                     (earlier milestone artifacts)
|   |-- walkthrough.md                           (initial baselines)
|   |-- proposal.tex, my_part.tex                (earlier LaTeX)
|   |-- jbd2_probe_module/                       (kprobe module source)
|   +-- jbd2_coalesce_module/                    (C1-era coalesce module)
|
+-- linux-6.1.4/                                 (gitignored - see below)
\end{lstlisting}

The Linux kernel source tree \texttt{linux-6.1.4/} is not in the
repository (too large). Reviewers should fetch the pristine 6.1.4
tarball from \url{https://cdn.kernel.org/pub/linux/kernel/v6.x/linux-6.1.4.tar.xz}
and apply the patch to it.

\subsection{Setup Instructions}

\paragraph{Hardware.}
\begin{itemize}[leftmargin=*]
  \item CPU: 2 cores minimum, 4+ cores recommended.
  \item Memory: 4~GB minimum, 8~GB recommended.
  \item Storage: 40~GB free for kernel build and tests. SSD or NVMe
        recommended for realistic fsync numbers; a VM with a loop
        device works but will produce smaller wall-time speedups
        (see main body for analysis).
  \item No GPU or other special hardware required.
\end{itemize}

A single-partition install is fine. The kernel source tree uses
about 15~GB during build, plus 500~MB each under \texttt{/boot} and
\texttt{/lib/modules/<version>} after installation.

\paragraph{Operating system.} Ubuntu 22.04 or 24.04 LTS on a VM or
bare-metal. We tested on Ubuntu~24.04.2 inside VirtualBox, Milan's
bare-metal was also Ubuntu-based.

\paragraph{Software dependencies.} Install once:

\begin{lstlisting}[language=bash]
sudo apt update
sudo apt install -y fio attr git build-essential libncurses-dev \
    bison flex libssl-dev libelf-dev bc dwarves zstd patch wget \
    gawk linux-tools-common linux-tools-$(uname -r)
\end{lstlisting}

For xfstests (optional, for detailed correctness evaluation):

\begin{lstlisting}[language=bash]
sudo apt install -y xfslibs-dev libattr1-dev libacl1-dev libaio-dev \
    acl libtool-bin e2fsprogs libcap-dev quota uuid-runtime xfsprogs \
    autoconf-archive libtool automake liburing-dev pkg-config \
    libgdbm-dev
\end{lstlisting}

\paragraph{Kernel build.}

\begin{lstlisting}[language=bash]
cd ~
wget https://cdn.kernel.org/pub/linux/kernel/v6.x/linux-6.1.4.tar.xz
tar xf linux-6.1.4.tar.xz
cd linux-6.1.4

# Apply the C3 patch
patch -p1 < ~/jbd2-project/our_optimization/fc-inline-xattr.patch

# Start from the running kernel's config, set a distinct LOCALVERSION
cp /boot/config-$(uname -r) .config
sed -i 's/^CONFIG_LOCALVERSION=.*/CONFIG_LOCALVERSION="-cs614-c3-patched"/' .config
make olddefconfig

# Build (20-45 min)
make -j$(nproc) bzImage modules

# Install (keeps existing kernels intact)
sudo rm -rf /lib/modules/6.1.4-cs614-c3-patched
sudo make modules_install
sudo make install
sudo update-grub

sudo reboot    # select 6.1.4-cs614-c3-patched in GRUB
\end{lstlisting}

After reboot, confirm with \texttt{uname -r}. Expected:
\texttt{6.1.4-cs614-c3-patched}.

\subsection{Features Implemented}

The patch implements one feature: Fast Commit support for inline
extended attribute changes in ext4 on Linux~6.1.4. Under the hood
this is two small changes in two files, covered in
Section~7.3 of this report.

For each supported scenario, the test script, parameters, and
expected outcome are below.

\begin{center}
\small
\begin{tabularx}{\textwidth}{lXXX}
\toprule
\# & Scenario & Script & Objective / Expected Outcome \\
\midrule
1 & Primary xattr workload. 5000
  \texttt{fsetxattr}+\texttt{fsync} ops; 256 MB loop fs;
  \texttt{-O fast\_commit}; single inline-sized value. &
  \texttt{bench\_xattr.sh} &
  Measure per-op latency and JBD2 commit count to confirm fast
  commit engages. Patched: wall time $\sim$23~s, tx\_delta
  $\sim$78. Stock: $\sim$31~s, tx\_delta 5000. 64$\times$ reduction
  in full commits; 27\% wall-time reduction. \\
\addlinespace
2 & Non-xattr control. fio 4-job randwrite+fsync for 30~s. &
  \texttt{bench\_async\_prefetch.sh} &
  Confirm unrelated workloads are not regressed. Patched IOPS
  within $\pm$5\% of stock. Observed: 528 vs 530. \\
\addlinespace
3 & Crash recovery A. 100 setfattr ops on one file; lazy-umount
  simulated crash; remount; count. &
  \texttt{c3\_crash\_test\_a.sh} &
  Expanded FC\_TAG\_INODE replay restores all inline xattrs. PASS
  with 100 / 100. \\
\addlinespace
4 & Crash recovery B. Set 100; remove 50 even-numbered;
  simulated crash. &
  \texttt{c3\_crash\_test\_b.sh} &
  Removal path replays correctly via fast commit. PASS with
  exactly 50 odd-numbered xattrs remaining. \\
\addlinespace
5 & Crash recovery C. 500-byte xattr (forces the block path,
  \texttt{touched\_block=true}); simulated crash. &
  \texttt{c3\_crash\_test\_c.sh} &
  Full-commit fallback is taken for block xattrs and data is
  preserved. PASS with all 668 bytes recovered. \\
\addlinespace
6 & xfstests xattr + fast commit subset. \texttt{generic/\{062,
  118, 300, 337, 454, 455, 473, 482\}} and \texttt{ext4/032}. &
  xfstests & 6/6 runnable tests pass identically on stock and
  patched. One pre-existing \texttt{generic/473} failure on both
  kernels (not our regression). \\
\bottomrule
\end{tabularx}
\end{center}

\paragraph{Findings during evaluation.} No crashes, deadlocks, or
assertion failures were observed during our evaluation of the
patched kernel. \texttt{dmesg} showed no ext4 or JBD2 WARN, ERROR,
BUG, or OOPS entries. The patched kernel mounts, operates, remounts,
and unmounts ext4 correctly.

Two environmental issues worth noting, both unrelated to the patch:
\begin{itemize}[leftmargin=*]
  \item \texttt{xfstests generic/062} uses \texttt{asort()} which
        is a gawk extension. On Ubuntu systems with the default
        \texttt{mawk}, the test reports an output mismatch that
        looks like a regression but is a test-framework issue.
        Install \texttt{gawk} to avoid this.
  \item \texttt{xfstests generic/388} is XFS-specific. Its
        \texttt{fsstress} phase hangs on ext4+fast\_commit
        regardless of our patch. We excluded it from our test list.
\end{itemize}

\subsection{Assumptions and Unsupported Features}

\paragraph{Covered.}
\begin{itemize}[leftmargin=*]
  \item Inline xattrs on ext4 (values that fit in the inode's
        extra region).
  \item SELinux labels, POSIX ACLs, short \texttt{user.*} xattrs.
  \item Both \texttt{setxattr} and \texttt{removexattr}.
  \item All ext4 journaling modes (\texttt{data=ordered}, etc.).
\end{itemize}

\paragraph{Not covered.} These paths are left on the full-commit
fallback by design:
\begin{itemize}[leftmargin=*]
  \item Xattrs in a separate 4~KB block (\texttt{i\_file\_acl} path).
        Large values go here. Still correct; just no speedup.
  \item Xattrs stored via the \texttt{ea\_inode} feature.
  \item Non-ext4 filesystems.
  \item First xattr on a fresh filesystem (by design; feature bit
        persistence).
\end{itemize}

No on-disk format changes, no compat bits, no new tag types, no
new mount options.

\subsection{Getting Started (within 30 minutes)}

This path verifies the artifact without rebuilding the kernel. It
exercises patch application, compile-check, and the raw result
data we recorded.

\paragraph{Step 1.} Clone the repository, or extract the submitted
ZIP. We assume the artifact is at \texttt{\textasciitilde/jbd2-project/}.

\paragraph{Step 2.} Verify the patch applies cleanly to pristine
Linux~6.1.4:

\begin{lstlisting}[language=bash]
cd ~
[ -d linux-6.1.4 ] || {
    wget -q https://cdn.kernel.org/pub/linux/kernel/v6.x/linux-6.1.4.tar.xz
    tar xf linux-6.1.4.tar.xz
}
patch -p1 --dry-run -d linux-6.1.4 \
      < ~/jbd2-project/our_optimization/fc-inline-xattr.patch
\end{lstlisting}

Expected output: \texttt{checking file fs/ext4/fast\_commit.c} and
\texttt{checking file fs/ext4/xattr.c} with no FAILED or malformed
hunks. Takes under a minute.

\paragraph{Step 3.} Compile-check the patched files in isolation
(no full kernel build):

\begin{lstlisting}[language=bash]
cd linux-6.1.4
patch -p1 < ~/jbd2-project/our_optimization/fc-inline-xattr.patch
cp /boot/config-$(uname -r) .config 2>/dev/null || make defconfig
make olddefconfig
make fs/ext4/xattr.o fs/ext4/fast_commit.o 2>&1 | tail -5
\end{lstlisting}

Expected: two \texttt{CC} lines for \texttt{xattr.o} and
\texttt{fast\_commit.o}, no \texttt{error:}. Takes 1-2~min on 4
cores.

\paragraph{Step 4.} Inspect recorded benchmark data:

\begin{lstlisting}[language=bash]
cd ~/jbd2-project
cat our_optimization/eval_results_c3/6.1.4-cs614-hacker/xattr_loop.txt
cat our_optimization/eval_results_c3/6.1.4-cs614-c3-patched/xattr_loop.txt
cat our_optimization/eval_results_c3/milan/STOCK_6.1.4/xattr_loop.txt
cat our_optimization/eval_results_c3/milan/PATCHED_6.1.4-C3Patch/xattr_loop.txt
\end{lstlisting}

These four files together give the VM and bare-metal comparison.

\paragraph{Step 5.} Read the full writeup:
\texttt{our\_optimization/CANDIDATE3\_results.md}.

\paragraph{Supplying your own inputs.} \texttt{bench\_xattr.sh}
is parameterized. To vary:
\begin{itemize}[leftmargin=*]
  \item \texttt{N}: iterations (default 5000).
  \item \texttt{IMG\_SIZE\_MB}: size of the loop-backed test
        filesystem.
  \item Value length: edit \texttt{xattr\_fsync\_helper.c} so
        \texttt{snprintf} produces a longer value. Beyond
        $\sim$80~bytes the xattr spills out of inline and you will
        see the full-commit fallback engage (\texttt{touched\_block
        = true}, \texttt{tx\_delta} returns toward \texttt{N}).
\end{itemize}

\subsection{Detailed Evaluation}

\paragraph{Experiment 1: Primary xattr speedup.}
\begin{itemize}[leftmargin=*]
  \item \textit{Purpose.} Quantify per-op latency and JBD2
        transaction-count reduction from our patch on a
        setxattr-heavy workload.
  \item \textit{How to run.} Boot into
        \texttt{6.1.4-cs614-c3-patched}, then
        \texttt{sudo bash bench\_xattr.sh}. Reboot into any
        non-C3 kernel (e.g., \texttt{6.1.4-cs614-hacker}), then
        re-run the same command.
  \item \textit{Estimated runtime.} 5~min patched + 5~min stock +
        two reboots.
  \item \textit{Expected result.} Patched elapsed
        22{,}000--25{,}000~ms; stock 30{,}000--35{,}000~ms.
        Patched \texttt{tx\_delta} under 100; stock
        \texttt{tx\_delta} $\approx$ \texttt{N}.
  \item \textit{How to access.} Each run writes
        \texttt{xattr\_loop.txt} in its result directory.
\end{itemize}

\paragraph{Experiment 2: Crash recovery correctness.}
\begin{itemize}[leftmargin=*]
  \item \textit{Purpose.} Verify our expanded FC\_TAG\_INODE replay
        restores inline xattrs, and that the full-commit fallback
        still works for block xattrs.
  \item \textit{How to run.} \texttt{sudo bash c3\_crash\_test\_a.sh}
        then \texttt{b.sh} then \texttt{c.sh}.
  \item \textit{Estimated runtime.} 30~seconds total.
  \item \textit{Expected result.} Each script prints
        \texttt{PASS: Test X}.
  \item \textit{How to access.} Output is the script's stdout.
\end{itemize}

\paragraph{Experiment 3: Non-xattr regression check.}
\begin{itemize}[leftmargin=*]
  \item \textit{Purpose.} Confirm the patch does not regress
        unrelated fsync workloads.
  \item \textit{How to run.}
        \texttt{sudo bash bench\_async\_prefetch.sh}.
  \item \textit{Estimated runtime.} 3-4~minutes.
  \item \textit{Expected result.} IOPS within $\pm$5\% of stock.
  \item \textit{How to access.} The script prints a fio summary
        and writes \texttt{eval\_results\_c2/<kernel>/*\_fio.txt}.
\end{itemize}

\paragraph{Experiment 4: xfstests subset.}
\begin{itemize}[leftmargin=*]
  \item \textit{Purpose.} Broader correctness validation against
        the upstream ext4 test suite.
  \item \textit{How to run.} Install xfstests (see dependencies),
        create the loop devices listed in
        \texttt{our\_optimization/xfstests\_c3/} setup notes, then:
        \texttt{sudo ./check generic/062 generic/118 generic/300
        generic/337 generic/454 generic/455 generic/473
        generic/482 ext4/032}.
  \item \textit{Estimated runtime.} 10-30~minutes.
  \item \textit{Expected result.} Same pass/fail pattern on stock
        and patched kernels: 6 pass, 3 not-applicable, 1
        pre-existing failure (\texttt{generic/473}).
  \item \textit{How to access.} \texttt{\textasciitilde/xfstests/results/}
        contains per-test output. Our recorded log is at
        \texttt{our\_optimization/xfstests\_c3/}.
\end{itemize}

\paragraph{Experiment 5: Bare-metal validation.}
\begin{itemize}[leftmargin=*]
  \item \textit{Purpose.} Confirm the VM result on real hardware.
  \item \textit{How to run.} Follow
        \texttt{INSTRUCTIONS\_MILAN\_C3.md} or
        \texttt{INSTRUCTIONS\_SAHIL\_C3.md} on a bare-metal
        Linux~6.1.4 machine.
  \item \textit{Estimated runtime.} 1-1.5~hours per machine
        including the kernel build.
  \item \textit{Expected result.} Larger wall-time speedup than VM
        (Milan's run gave 48\%); same $\sim$64$\times$ transaction
        reduction.
  \item \textit{How to access.}
        \texttt{our\_optimization/eval\_results\_c3/<contributor>/*/xattr\_loop.txt}
        after the contributor's branch is merged.
\end{itemize}

\end{document}
